\documentclass[11pt]{article}
\usepackage{amsmath, amssymb}
\usepackage[margin=1in]{geometry}
\usepackage{graphicx}   % for \resizebox
\usepackage{xcolor}

% --- notation (hats over the letters only) ---
\newcommand{\ITTY}{\mathrm{ITT}_Y}
\newcommand{\ITTD}{\mathrm{ITT}_D}
\newcommand{\CACE}{\mathrm{CACE}}
\newcommand{\hITTY}{\widehat{\mathrm{ITT}}_Y}
\newcommand{\hITTD}{\widehat{\mathrm{ITT}}_D}
\newcommand{\hCACE}{\widehat{\mathrm{CACE}}}
\newcommand{\E}{\mathbb{E}}

\begin{document}

\begin{center}
{\large\bfseries Why $\dfrac{(\hCACE-\CACE)-L}{\sqrt{V_\Delta}}\overset{p}{\longrightarrow}0$:
an exact look at the remainder}
\end{center}

\bigskip
\noindent
The claim on the slide is that the two standardized estimators differ by a
remainder that vanishes:
\begin{equation*}
\resizebox{\linewidth}{!}{$\displaystyle
\underbrace{\frac{\hCACE - \CACE}{\sqrt{V_\Delta}}}_{\text{(a) estimator}}
- \underbrace{\frac{L}{\sqrt{V_\Delta}}}_{\text{(b) approximation}}
= \frac{\left(\hCACE - \CACE\right) - L}{\sqrt{V_\Delta}}
\overset{p}{\longrightarrow} 0$}
\end{equation*}
We will compute the numerator $\left(\hCACE-\CACE\right)-L$ \emph{exactly} and
see precisely why it is negligible.

\bigskip
\noindent\textbf{Definitions.}
\begin{equation*}
\hCACE = \frac{\hITTY}{\hITTD}, \qquad
\CACE = \frac{\ITTY}{\ITTD}, \qquad
L = \frac{\hITTY - \CACE\,\hITTD}{\ITTD}.
\end{equation*}
Write the estimation errors of the two ingredients as
\begin{equation*}
\Delta_Y = \hITTY - \ITTY, \qquad \Delta_D = \hITTD - \ITTD,
\end{equation*}
so that $\hITTY = \ITTY + \Delta_Y$ and $\hITTD = \ITTD + \Delta_D$.

\bigskip
\noindent\textbf{Step 1. The exact estimation error, over a common denominator.}
\begin{align*}
\hCACE - \CACE
&= \frac{\hITTY}{\hITTD} - \frac{\ITTY}{\ITTD}
= \frac{\ITTD\,\hITTY - \ITTY\,\hITTD}{\ITTD\,\hITTD} \\[4pt]
&= \frac{\ITTD\left(\ITTY+\Delta_Y\right) - \ITTY\left(\ITTD+\Delta_D\right)}{\ITTD\,\hITTD}
= \frac{\ITTD\,\Delta_Y - \ITTY\,\Delta_D}{\ITTD\,\hITTD}.
\end{align*}

\bigskip
\noindent\textbf{Step 2. The linear approximation, in the same notation.}
Since $\CACE = \ITTY/\ITTD$,
\begin{align*}
L = \frac{\hITTY - \CACE\,\hITTD}{\ITTD}
&= \frac{\hITTY - \frac{\ITTY}{\ITTD}\,\hITTD}{\ITTD}
= \frac{\ITTD\,\hITTY - \ITTY\,\hITTD}{\ITTD^{2}} \\[4pt]
&= \frac{\ITTD\,\Delta_Y - \ITTY\,\Delta_D}{\ITTD^{2}}.
\end{align*}

\bigskip
\noindent\textbf{Key observation.}
The estimator's error and its linear approximation have the \emph{same numerator}
$\;\ITTD\,\Delta_Y - \ITTY\,\Delta_D\;$; they differ \emph{only} in the
denominator ($\ITTD\,\hITTD$ versus $\ITTD^{2}$). Hence
\begin{equation*}
\hCACE - \CACE
= L \cdot \frac{\ITTD^{2}}{\ITTD\,\hITTD}
= L \cdot \frac{\ITTD}{\hITTD}.
\end{equation*}

\bigskip
\noindent\textbf{Step 3. The remainder.}
\begin{align*}
\left(\hCACE - \CACE\right) - L
&= L\,\frac{\ITTD}{\hITTD} - L
= L\left(\frac{\ITTD}{\hITTD} - 1\right)
= L\,\frac{\ITTD - \hITTD}{\hITTD} \\[4pt]
&= -\,L\,\frac{\hITTD - \ITTD}{\hITTD}
= -\,L\,\frac{\Delta_D}{\hITTD}.
\end{align*}
That is,
\begin{equation*}
\boxed{\;\left(\hCACE - \CACE\right) - L
= -\,L\,\frac{\hITTD - \ITTD}{\hITTD}\;}
\end{equation*}
the linear approximation $L$ times (minus) the \emph{relative error of the first
stage}, $\Delta_D/\hITTD = (\hITTD - \ITTD)/\hITTD$.

\bigskip
\noindent\textbf{Step 4. Standardize: the remainder is a product of two factors.}
Dividing the boxed identity by $\sqrt{V_\Delta}$,
\begin{equation*}
\frac{\left(\hCACE - \CACE\right) - L}{\sqrt{V_\Delta}}
= -\,\underbrace{\frac{L}{\sqrt{V_\Delta}}}_{\textstyle \text{Factor 1}}
\;\cdot\;
\underbrace{\frac{\hITTD - \ITTD}{\hITTD}}_{\textstyle \text{Factor 2}} .
\end{equation*}
We must show this product converges in probability to $0$. The plan: \emph{Factor 1
stays bounded}, \emph{Factor 2 vanishes}, and a bounded quantity times a vanishing
one must vanish. Throughout, ``$N$'' indexes a growing sequence of finite
populations; the only randomness is the assignment.

\bigskip
\noindent\textbf{The rates, in one breath (the intuition).}
A difference-in-means has sampling variance equal to a sum of within-arm variances
divided by the arm sizes, and the arm sizes grow in proportion to $N$. So its
\emph{variance shrinks in proportion to $1/N$}, and its \emph{standard deviation in
proportion to $1/\sqrt{N}$}. Hence:
\begin{itemize}
\item A typical first-stage error $\Delta_D = \hITTD - \ITTD$ is about
$1/\sqrt{N}$ in size and shrinks to $0$ as the study grows.
\item $L$ and $\sqrt{V_\Delta}$ are \emph{both} about $1/\sqrt{N}$: $L$ is a linear
combination of the two errors, and $\sqrt{V_\Delta}$ is exactly $L$'s own standard
deviation. Their ratio (Factor 1) is therefore about $1$ --- it neither grows nor
shrinks, which is precisely why it can have a non-degenerate ($N(0,1)$) limit.
\item The denominator $\hITTD$ settles on the fixed nonzero number $\ITTD$, so
Factor 2 is about $(1/\sqrt{N})/\ITTD$ --- it shrinks to $0$ at the $1/\sqrt{N}$
rate.
\end{itemize}
So the product is about $1 \times 1/\sqrt{N} = 1/\sqrt{N}$: the standardized
remainder shrinks to $0$ at the root-$N$ rate. The rest makes this rigorous.

\bigskip
\noindent\textbf{Factor 2 vanishes (consistency of the first stage).}
``$W_N$ converges in probability to $0$'' means: \emph{for every tolerance
$\varepsilon > 0$, the probability $\Pr(\lvert W_N\rvert > \varepsilon)$ can be made
as small as we like by taking $N$ large enough.} The first-stage estimator is
unbiased, $\E[\hITTD] = \ITTD$, with variance proportional to $1/N$. Chebyshev's
inequality then gives, for any fixed $\varepsilon>0$,
\begin{equation*}
\Pr\!\left(\lvert \Delta_D\rvert > \varepsilon\right)
\;\le\; \frac{\operatorname{Var}(\hITTD)}{\varepsilon^{2}}
\;\propto\; \frac{1}{N\,\varepsilon^{2}} \;\longrightarrow\; 0 ,
\end{equation*}
so $\Delta_D \to 0$ in probability. Because $\ITTD \neq 0$ and
$\hITTD = \ITTD + \Delta_D$ settles on $\ITTD$, the ratio Factor 2 has numerator
going to $0$ and denominator going to a nonzero constant, so
\begin{equation*}
\text{Factor 2} = \frac{\Delta_D}{\hITTD} \;\longrightarrow\; 0
\quad\text{in probability.}
\end{equation*}

\bigskip
\noindent\textbf{Factor 1 stays bounded (it cannot blow up).}
``$X_N$ is bounded in probability'' means: \emph{the values do not escape to
infinity --- for every small $\delta>0$ there is a finite cutoff $M$, the same for
all $N$, with $\Pr(\lvert X_N\rvert > M) < \delta$.} Factor 1 $=L/\sqrt{V_\Delta}$
converges in distribution to $N(0,1)$ (a linear combination of jointly Normal
differences in means, divided by its own standard deviation). A sequence whose
distribution settles on a fixed, genuine distribution cannot leak mass off to
$\pm\infty$: for large $N$ its distribution is close to $N(0,1)$, and the standard
Normal puts all but $\delta$ of its mass inside $[-M, M]$ once $M$ is the
$(1-\delta/2)$ Normal quantile. So $\Pr(\lvert \text{Factor 1}\rvert > M)$ is
eventually below $\delta$: Factor 1 does \emph{not} blow up.

\bigskip
\noindent\textbf{Bounded $\times$ vanishing $=$ vanishing (the one lemma we need).}
Let $X_N$ be bounded in probability and $Y_N \to 0$ in probability. Then
$X_N Y_N \to 0$ in probability. \emph{Proof.} Fix $\varepsilon>0$ and $\delta>0$.
Choose $M$ (finite, the same for all $N$) with $\Pr(\lvert X_N\rvert > M) <
\delta/2$. Whenever $\lvert X_N\rvert \le M$ we have $\lvert X_N Y_N\rvert \le
M\,\lvert Y_N\rvert$, so $\lvert X_N Y_N\rvert > \varepsilon$ forces
$\lvert Y_N\rvert > \varepsilon/M$. Hence
\begin{equation*}
\Pr\!\left(\lvert X_N Y_N\rvert > \varepsilon\right)
\;\le\; \Pr\!\left(\lvert X_N\rvert > M\right)
+ \Pr\!\left(\lvert Y_N\rvert > \tfrac{\varepsilon}{M}\right)
\;<\; \tfrac{\delta}{2} + \Pr\!\left(\lvert Y_N\rvert > \tfrac{\varepsilon}{M}\right).
\end{equation*}
Since $\varepsilon/M$ is a fixed positive number and $Y_N \to 0$ in probability, the
last term is below $\delta/2$ for $N$ large, so the bound is below $\delta$. As
$\varepsilon,\delta$ were arbitrary, $X_N Y_N \to 0$ in probability. \hfill$\square$

\bigskip
\noindent\textbf{Conclusion.}
With $X_N = $ Factor 1 (bounded) and $Y_N = $ Factor 2 (vanishing), the lemma gives
\begin{equation*}
\frac{\left(\hCACE - \CACE\right) - L}{\sqrt{V_\Delta}}
= -\,X_N\,Y_N \;\longrightarrow\; 0 \quad\text{in probability.}
\end{equation*}
And since Factor 1 $=L/\sqrt{V_\Delta} \to N(0,1)$ while the standardized remainder
$\to 0$ in probability, adding them back,
\begin{equation*}
\frac{\hCACE - \CACE}{\sqrt{V_\Delta}}
= \frac{L}{\sqrt{V_\Delta}} + \frac{\left(\hCACE - \CACE\right) - L}{\sqrt{V_\Delta}}
\;\longrightarrow\; N(0,1),
\end{equation*}
by Slutsky's lemma (a sequence converging in distribution, plus one vanishing in
probability, keeps the same limiting distribution). \hfill$\blacksquare$

\bigskip
\noindent\textbf{The ``second order'' story, in one line.}
Substituting $L = (\ITTD\,\Delta_Y - \ITTY\,\Delta_D)/\ITTD^{2}$ into the boxed
identity,
\begin{equation*}
\left(\hCACE - \CACE\right) - L
= -\,\frac{\left(\ITTD\,\Delta_Y - \ITTY\,\Delta_D\right)\,\Delta_D}
{\ITTD^{2}\,\hITTD},
\end{equation*}
a product of \emph{two} errors ($\Delta_Y\Delta_D$ and $\Delta_D^{2}$). Each error
is about $1/\sqrt{N}$, so their product is about $1/N$; dividing by
$\sqrt{V_\Delta}\approx 1/\sqrt{N}$ leaves about $1/\sqrt{N}\to 0$. The estimator's
error $L$ is \emph{first order} (about $1/\sqrt{N}$); the gap the tangent misses is
\emph{second order} (about $1/N$) --- smaller by a factor of about $1/\sqrt{N}$.
That is the precise sense in which the remainder is negligible relative to the
standard error.

\end{document}
